\documentclass{article}

% packages
\usepackage{amsmath, amsthm, thmtools, amsfonts, amssymb, luacode, catchfile, tikzducks, hyperref, ifthen}
\ifcsname c@kobocompile\endcsname
	\usepackage[a5paper, total={1072pt, 1448pt}, margin=10pt, includeheadfoot]{geometry} % set page margins
\else
	\usepackage[a4paper, margin=50pt, includeheadfoot]{geometry}
\fi
\usepackage[shortlabels]{enumitem}
\usepackage[skip=3pt, indent=0pt]{parskip}

% language
\usepackage[bidi=basic, layout=tabular, provide=*]{babel}
\ifcsname c@english\endcsname
	\babelprovide[main, import]{english}
\else
	\babelprovide[main, import]{hebrew}
	\babelprovide{rl}
\fi
%\babelfont{rm}{Libertinus Serif}
\babelfont{rm}[Renderer=Harfbuzz]{Libertinus Serif}
\babelfont{sf}{Libertinus Sans}
\babelfont{tt}{Libertinus Mono}

% style
\AddToHook{cmd/section/before}{\clearpage}	% Add line break before section
\linespread{1.3}
\setcounter{secnumdepth}{0}		% Remove default number tags from sections, this won't do well with theorems
\AtBeginDocument{\setlength{\belowdisplayskip}{3pt}}
\AtBeginDocument{\setlength{\abovedisplayskip}{3pt}}
\graphicspath{ {../images/} }

% operators
\DeclareMathOperator\cis{cis}
\DeclareMathOperator\Sp{Sp}
\DeclareMathOperator\tr{tr}
\DeclareMathOperator\im{Im}
\DeclareMathOperator\re{Re}
\DeclareMathOperator\diag{diag}
\DeclareMathOperator*\lowlim{\underline{lim}}
\DeclareMathOperator*\uplim{\overline{lim}}
\DeclareMathOperator\rng{rng}
\DeclareMathOperator\Sym{Sym}
\DeclareMathOperator\Arg{Arg}
\DeclareMathOperator\Log{Log}
\DeclareMathOperator\dom{dom}
\DeclareMathOperator\supp{Supp}
\DeclareMathOperator\var{Var}
\DeclareMathOperator\cov{Cov}

% commands
%\renewcommand\qedsymbol{\textbf{מש''ל}}
%\renewcommand\qedsymbol{\fbox{\emoji{lizard}}}
\newcommand{\Aa}[0]{\mathcal{A}}
\newcommand{\Bb}[0]{\mathcal{B}}
\newcommand{\CC}[0]{\mathbb{C}}
\newcommand{\Cc}[0]{\mathcal{C}}
\newcommand{\EE}[0]{\mathbb{E}}
\newcommand{\FF}[0]{\mathbb{F}}
\newcommand{\Ff}[0]{\mathcal{F}}
\newcommand{\Ii}[0]{\mathcal{I}}
\newcommand{\Gg}[0]{\mathcal{G}}
\newcommand{\Ll}[0]{\mathcal{L}}
\newcommand{\Mm}[0]{\mathcal{M}}
\newcommand{\NN}[0]{\mathbb{N}}
\newcommand{\Nn}[0]{\mathcal{N}}
\newcommand{\PP}[0]{\mathbb{P}}
\newcommand{\Pp}[0]{\mathcal{P}}
\newcommand{\QQ}[0]{\mathbb{Q}}
\newcommand{\RR}[0]{\mathbb{R}}
\newcommand{\Rr}[0]{\mathcal{R}}
\newcommand{\Ss}[0]{\mathcal{S}}
\newcommand{\TT}[0]{\mathbb{T}}
\newcommand{\Uu}[0]{\mathcal{U}}
\newcommand{\Vv}[0]{\mathcal{V}}
\newcommand{\Ww}[0]{\mathcal{W}}
\newcommand{\ZZ}[0]{\mathbb{Z}}
\newcommand{\acts}[0]{\circlearrowright}
\newcommand{\explain}[2] {
	\begin{flalign*}
		 && \text{#2} && \text{#1}
	\end{flalign*}
}
\newcommand{\maketitleprint}[0]{ \begin{center}
	%\begin{tikzpicture}[scale=3]
	%	\duck[graduate=gray!20!black, tassel=red!70!black]
	%\end{tikzpicture}	
	\includegraphics[width=6cm]{cover}
\end{center}
}

% theorem commands
\newtheoremstyle{c_remark}
	{}	% Space above
	{}	% Space below
	{}% Body font
	{}	% Indent amount
	{\bfseries}	% Theorem head font
	{}	% Punctuation after theorem head
	{.5em}	% Space after theorem head
	{\thmname{#1}\thmnumber{ #2}\thmnote{ \normalfont{\text{(#3)}}}}	% head content
\newtheoremstyle{c_definition}
	{3pt}	% Space above
	{3pt}	% Space below
	{}% Body font
	{}	% Indent amount
	{\bfseries}	% Theorem head font
	{}	% Punctuation after theorem head
	{.5em}	% Space after theorem head
	{\thmname{#1}\thmnumber{ #2}\thmnote{ \normalfont{\text{(#3)}}}}	% head content
\newtheoremstyle{c_plain}
	{3pt}	% Space above
	{3pt}	% Space below
	{\itshape}% Body font
	{}	% Indent amount
	{\bfseries}	% Theorem head font
	{}	% Punctuation after theorem head
	{.5em}	% Space after theorem head
	{\thmname{#1}\thmnumber{ #2}\thmnote{ \text{(#3)}}}	% head content

\ifcsname c@english\endcsname
	\theoremstyle{plain}
	\newtheorem{theorem}{Theorem}[section]
	\newtheorem{lemma}[theorem]{Lemma}
	\newtheorem{proposition}[theorem]{Proposition}
	\newtheorem*{proposition*}{Proposition}
	%\newtheorem{corollary}[theorem]{אין חלופה עברית}

	\theoremstyle{definition}
	\newtheorem{definition}[theorem]{Definition}
	\newtheorem*{definition*}{Definition}
	\newtheorem{example}{Example}[section]
	\newtheorem{exercise}{Exercise}[section]

	\theoremstyle{remark}
	\newtheorem*{remark}{Remark}
	\newtheorem*{solution}{Solution}
	\newtheorem{conclusion}[theorem]{Conclusion}
	\newtheorem{notation}[theorem]{Notation}
\else
	\theoremstyle{c_plain}
	\newtheorem{theorem}{משפט}[section]
	\newtheorem{lemma}[theorem]{למה}
	\newtheorem{proposition}[theorem]{טענה}
	\newtheorem*{proposition*}{טענה}
	%\newtheorem{corollary}[theorem]{אין חלופה עברית}

	\theoremstyle{c_definition}
	\newtheorem{definition}[theorem]{הגדרה}
	\newtheorem*{definition*}{הגדרה}
	\newtheorem{example}{דוגמה}[section]
	\newtheorem{exercise}{תרגיל}[section]

	\theoremstyle{c_remark}
	\newtheorem*{remark}{הערה}
	\newtheorem*{solution}{פתרון}
	\newtheorem{conclusion}[theorem]{מסקנה}
	\newtheorem{notation}[theorem]{סימון}
\fi

% Questions related commands
\newcounter{question}
\setcounter{question}{1}
\newcounter{sub_question}
\setcounter{sub_question}{1}

\ifcsname c@english\endcsname
	\newcommand{\question}[1][0]{
		\ifthenelse{#1 = 0}{}{\setcounter{question}{#1}}
		\section{Question \arabic{question}}
		\addtocounter{question}{1}
		\setcounter{sub_question}{1}
	}

	\newcommand{\subquestion}[1][0]{
		\ifthenelse{#1 = 0}{}{\setcounter{sub_question}{#1}}
		\subsection{Part \alph{sub_question}}
		\addtocounter{sub_question}{1}
	}
\else
	\newcommand{\question}[1][0]{
		\ifthenelse{#1 = 0}{}{\setcounter{question}{#1}}
		\section{שאלה \arabic{question}}
		\addtocounter{question}{1}
		\setcounter{sub_question}{1}
	}

	\newcommand{\subquestion}[1][0]{
		\ifthenelse{#1 = 0}{}{\setcounter{sub_question}{#1}}
		\subsection{סעיף \localecounter{letters.gershayim}{sub_question}}
		\addtocounter{sub_question}{1}
	}
\fi

% import lua and start of document
\directlua{common = require ('../common')}

\GetEnv{AUTHOR}

% headers
\author{\AUTHOR}
\date\today

\title{פתרון ממ''ן 13 – חשבון אינפיניטסימלי 1 (20474)}

\begin{document}
\maketitle
\section{שאלה 1}
\subsection{סעיף א'}
נגדיר $a_1 = 0$ ולכל $n$:
\[
	a_{n + 1} = \frac{1}{4(1 - a_n)}
\]
נוכיח כי הסדרה מוגדרת לכל $n$
\begin{proof}
	איבר איננו מוגדר כאשר $a_n = 1$, במקרה זה האיבר שלאחריו $a_{n + 1}$ לא מוגדר ולכן גם הסדרה לא תהיה מוגדרת לכל $n \in \NN$. \\*
	נוכיח באינדוקציה שמקרה זה לא מתקיים על־ידי ההוכחה כי $0 < a_n < \frac{1}{2}$ לכל $n > 1$. \\*
	\textbf{בסיס האינדוקציה:} על־פי חישוב $a_2 = \frac{1}{4(1 - 0)} = \frac{1}{4}$ ולכן $0 < a_2 < \frac{1}{2}$. \\*
	\textbf{מהלך האינדוקציה:} נניח כי $0 < a_n < \frac{1}{2}$ ונראה כי $0 < a_{n + 1} < \frac{1}{2}$. \\*
	מתקיים
	\begin{align*}
		& 0 < a_n < \frac{1}{2} \\
		& 1 - 0 = 1 > 1 - a_n > \frac{1}{2} = 1 - \frac{1}{2} \\
		& 4 > 4(1 - a_n) > 2 \\
		& 0 < \frac{1}{4} < \frac{1}{4(1 - a_n)} < \frac{1}{2} \\
		& 0 < a_{n + 1} < \frac{1}{2} \\
	\end{align*}
	מהלך האינדוקציה הושלם ולכן לכל $n > 1$ מתקיים
	\[
		0 < a_n < \frac{1}{2}
	\]
	לפיכך אנו יודעים גם כי לכל $n > 1$ מתקיים $a_n \ne 1$ ולכן $a_n$ מוגדר לכל $n$.
\end{proof}

\subsection{סעיף ב'}
נוכיח כי הסדרה $(a_n)$ מתכנסת ונמצא את ערך גבולה.
\begin{proof}
	תחילה נוכיח באינדוקציה כי לכל $n > 1$ מתקיים $a_{n + 1} > a_n$: \\*
	\textbf{בסיס האינדוקציה:}
	\[
		a_3 = \frac{1}{4(1 - a_2)} = \frac{1}{4 - 1} = \frac{1}{3} > \frac{1}{4} = a_2
	\]
	\textbf{מהלך האינדוקציה:}
	נניח כי $a_n > a_{n - 1}$ ונוכיח כי $a_{n + 1} > a_n$.
	\begin{align*}
		& a_n > a_{n - 1} \\
		& 1 - a_n < 1 - a_{n - 1} \\
		& 4(1 - a_n) < 4(1 - a_{n - 1}) \\
		& \frac{1}{4(1 - a_n)} > \frac{1}{4(1 - a_{n - 1})} \\
		& a_{n + 1} > a_n
	\end{align*}
	התנאי מתקיים ולכן הסדרה עולה לכל $n > 1$. \\*
	ידוע כי $a_n < \frac{1}{2}$ לכל $n$, וראינו כי הסדרה עולה ולכן הסדרה $(a_n)$ חסומה ועולה, ולכן לפי משפט 3.16 מתכנסת.
	נגדיר $\lim_{n \to \infty} a_n = L$. \\*
	על־פי אריתמטיקה של הגבולות:
	\begin{align*}
		L & = \lim_{n \to \infty} a_{n + 1} \\
		  & = \lim_{n \to \infty} \frac{1}{4(1 - a_n)}
		  & = \frac{1}{\displaystyle\lim_{n \to \infty} 4(1 - a_n)} \\
		  & = \frac{1}{4 \displaystyle\lim_{n \to \infty} (1 - a_n)}
		  & = \frac{1}{4 (1 - \displaystyle\lim_{n \to \infty} a_n)} \\
		L & = \frac{1}{4 (1 - L)}
		  & \rightarrow 4L(1 - L) = 1
	\end{align*}
	נחשב את ערכו של $L$:
	\begin{align*}
		& 4L - 4L^2 - 1 = 0 \\
		& 4L^2 - 4L + 1 = 0 \\
		& L = \frac{4 \pm \sqrt{4^2 - 4 \cdot 4}}{2 \cdot 4} = \frac{4 \pm 0}{8} = \frac{1}{2}
	\end{align*}
	לכן לסדרה $(a_n)$ יש גבול והוא
	\[
		\lim_{n \to \infty} a_n = \frac{1}{2}
	\]
\end{proof}

\section{שאלה 2}
\subsection{סעיף א'}
נחשב את הגבול
\[
	\lim_{n \to \infty}
	\frac{{(-5)}^n + 2{(-2)}^n + 3}{5^{n + 1} + 2{(-3)}^n + 3}
\]
נגדיר שתי תת־סדרות $(a_{n_k})$ ו־$(a_{m_k})$
המוגדרות כסדרת האיברים הזוגיים והאי־זוגיים בהתאמה של הסדרה.
סדרות אלה מכסות את הסדרה המקורית, ומתקיים:
\[
	a_{n_k} = \frac{{5}^n + 2 \cdot 2^n + 3}{5^{n + 1} + 2 \cdot 3^n + 3}
\]
נחשב את גבולה בעזרת אריתמטיקה של גבולות:
\begin{align*}
	\lim_{n \to \infty} a_{n_k}
	& = \lim_{n \to \infty} \frac{{5}^n + 2 \cdot 2^n + 3}{5^{n + 1} + 2 \cdot 3^n + 3} \\
	& = \lim_{n \to \infty} \frac{{5}^n/5^n + 2 \cdot 2^n / 5^n + 3/5^n}{5^{n + 1}/5^n + 2 \cdot 3^n/5^n + 3/5^n} \\
	& = \lim_{n \to \infty} \frac{1 + 0 + 0}{5^1 + 0 + 0} \\
	\lim_{n \to \infty} a_{n_k} & = \frac{1}{5}
\end{align*}
באופן דומה נראה כי
\[
	a_{m_k} = \frac{-{5}^m - 2 \cdot 2^m + 3}{5^{m + 1} - 2 \cdot 3^m + 3}
\]
חישוב דומה יוביל אותנו למסקנה
\[
	\lim_{m \to \infty} (a_{m_k}) = -\frac{1}{5}
\]
ראינו כי גבולות תת־הסדרות של הסדרה מתכנסות לערכים שונים,
לכן לפי משפט 3.31 הסדרה עצמה מתבדרת,
וגבולותיה החלקיים הם $-\frac{1}{5}, \frac{1}{5}$.

\subsection{סעיף ב'}
נחשב את הגבול
\[
	\lim_{n \to \infty} \frac{{(-5)}^n + 4^{n + 1} + 3}{{({-4})}^n + 2{(-2)}^n + 3}
\]
נגדיר סדרות זוגיות ואי־זוגיות כבסעיף א' ונחשב את גבולן:
\[
	\lim_{n \to \infty} \frac{5^n + 4^{n + 1} + 3}{4^n + 2 \cdot 2^n + 3}
	= \lim_{n \to \infty} \frac{5^n/5^n + 4^{n + 1}/5^n + 3/5^n}{4^n/5^n + 2 \cdot 2^n/5^n + 3/5^n}
	= \lim_{n \to \infty} \frac{1}{0^+}
	= \infty
\]
גבול האברים האי־זוגיים הוא, על־פי חישוב דומה:
\[
	\lim_{m \to \infty} \frac{-5^m + 4^{m + 1} + 3}{-4^m - 2 \cdot 2^m + 3}
	= \lim_{m \to \infty} \frac{-1}{0^-}
	= \infty
\]
על־פי משפט 3.31 מתקיים:
\[
	\lim_{n \to \infty} \frac{{(-5)}^n + 4^{n + 1} + 3}{{({-4})}^n + 2{(-2)}^n + 3} = \infty
\]

\subsection{סעיף ג'}
נוכיח כי לא מתקיים הגבול
\[
	\lim_{n \to \infty} {\left( \frac{1}{n} - 1 \right)}^n
\]
נגדיר $(a_{n_k})$ סדרת האיברים הזוגיים של $(a_n)$, $(a_{m_k})$ סדרת האיברים האי־זוגיים של $(a_n)$. \\*
נשים לב כי בסדרה $(a_{n_k})$ מתקיים
\[
	{\left( \frac{1}{n} - 1 \right)}^n
	= {(-1)}^n {\left( 1 - \frac{1}{n} \right)}^n
	= {\left( 1 - \frac{1}{n} \right)}^n
	\tag{1}
\]
כמו־כן מתקיים
\[
	{\left( 1 - \frac{1}{n} \right)}^n
	= {\left( \frac{n - 1}{n} \right)}^n
	= {\left( \frac{n - 1 + 1}{n - 1} \right)}^{-n}
	= {\left( 1 + \frac{1}{n - 1} \right)}^{-n}
\]
על־פי דוגמה 3.5 ושאלה 20 סעיף א' מתקיים:
\[
	\lim_{n \to \infty} a_{n_k}
	= \lim_{n \to \infty} {\left( 1 + \frac{1}{n - 1} \right)}^{-n}
	= e^{-1}
\]
באופן דומה עבור $(a_{m_k})$ ו־$(1)$ מתקיים
\[
	{\left( \frac{1}{n} - 1 \right)}^n
	= -{\left( 1 - \frac{1}{n} \right)}^n
\]
ולכן על־פי אריתמטיקה של גבולות
\[
	\lim_{m \to \infty} a_{m_k} = - a_{n_k} = -e^{-1}
\]
אנו רואים כי $(a_n)$ לא מתכנסת וגבולותיה החלקיים הם $\pm e^{-1}$.

\subsection{סעיף ד'}
תהי $(a_n)$ סדרה עולה ממש של מספרים שלמים, נוכיח כי מתקיים הגבול
\[
	\lim_{n \to \infty} {\left( 1 + \frac{1}{a_n} \right)}^{a_n}
\]
על־פי הגדרה 3.24 ומשפט 3.25 מתקיים
\[
	\lim_{n \to \infty} {\left( 1 + \frac{1}{a_n} \right)}^{a_n}
	= \lim_{n \to \infty} {\left( 1 + \frac{1}{n} \right)}^{n}
\]
ולכן על־פי דוגמה 3.5 הגבול מתקיים וערכו הוא $e$.

\section{שאלה 3}
תהי $a_n = \langle \sqrt{n} \rangle$.
\subsection{סעיף א'}
נוכיח כי הסדרה $(a_n)$ חסומה.
\begin{proof}
	נגדיר $l = \sqrt{n}$, אז מתקיים $a_n = \langle l \rangle$. על־פי הגדרת החלק השברי מתקיים $0 \le l < 1$, לכן מתקיים גם $0 \le \sqrt{n} < 1$. \\*
	נעלה בריבוע ונקבל $0^2 = 0 \le n < 1^2 = 1$. אנו רואים כי לכל $n \in \NN$ מתקיים $0 \le a_n < 1$ ולכן הסדרה $(a_n)$ חסומה ב־1.
\end{proof}

\subsection{סעיף ב'}
נחשב את $\displaystyle\lowlim_{n \to \infty} a_n$. \\*
נגדיר $n_k = n^2$, נשים לב כי מתקיים $a_{n_k} = \langle \sqrt{n^2} \rangle = \langle n \rangle = 0$.
לכן גם מתקיים $\lim_{n \to \infty} a_{n_k} = 0$ ובהתאם 0 הוא גבול חלקי של $(a_n)$. \\*
בסעיף הקודם הוכחנו שלכל $n$ מתקיים $0 \le a_n$ ולכן לא יתכן שקיים גבול הקטן מ־0, ובהתאם 0 הוא הגבול החלקי הקטן ביותר של $(a_n)$ ומתקיים
\[
	\lowlim_{n \to \infty} a_n = 0
\]

\subsection{סעיף ג'}
נגדיר $A = \{ a_n \mid n \in \NN \}$. נמצא את $\inf A$. \\*
כפי שראינו בסעיף הקודם, $0 \in A$ מקיים $\forall a \in A, 0 \le a$ ולכן מהגדרה 3.12 נובע ש־$\inf A = 0$. \\*
לקבוצה כמובן יש מינימום לפי טענה 3.13, הוא $0$.

\subsection{סעיף ד'}
נוכיח כי לכל $n \in \NN$ מתקיים
\[
	\left\langle \sqrt{n^2 - 1} \right\rangle = \sqrt{n^2 - 1} - n + 1
\]
\begin{proof}
	נראה כי
	\begin{align*}
		& -n \le -1 \\
		& -2n \le -2 \\
		& -2n + 1 \le -1 \\
		& n^2 - 2n + 1 \le n^2 - 1 \le n^2 & \text{מתקיים גם} \\
		& {(n - 1)}^2 \le n^2 - 1 \le n^2 \\
		& n - 1 \le \sqrt{n^2 - 1} \le n \\
		& \left\lfloor \sqrt{n^2 - 1} \right\rfloor = n - 1 & \text{לפי טענה 1.64 3.} \\
		& \sqrt{n^2 - 1} - \left\lfloor \sqrt{n^2 - 1} \right\rfloor = \left\langle \sqrt{n^2 - 1} \right\rangle = \sqrt{n^2 - 1} - n + 1
	\end{align*}
\end{proof}

\subsection{סעיף ה'}
נגדיר $(b_n)$ סדרה כך שמתקיים
\[
	b_n = \sqrt{n^2 - 1} - n + 1
\]
נוכיח כי יש גבול לסדרה, נשתמש באריתמטיקה של הגבולות, אילו קיים גבול אז יוצדקו מעברים אלה.
\begin{align*}
	\lim_{n \to \infty} \sqrt{n^2 - 1} - n + 1
	& = \lim_{n \to \infty}\sqrt{n + 1} \sqrt{n - 1} - \sqrt{n - 1} \sqrt{n - 1} \\
	& = \lim_{n \to \infty}\sqrt{n - 1} \left( \sqrt{n + 1} - \sqrt{n - 1} \right) \\
	& = \lim_{n \to \infty}\sqrt{n - 1} \frac{\left( \sqrt{n + 1} - \sqrt{n - 1} \right)\left( \sqrt{n + 1} + \sqrt{n - 1} \right)}
		{\left( \sqrt{n + 1} + \sqrt{n - 1} \right)} \\
	& = \lim_{n \to \infty}\sqrt{n - 1} \frac{n + 1 - n + 1}{\left( \sqrt{n + 1} + \sqrt{n - 1} \right)} \\
	& = \lim_{n \to \infty}\frac{2 \sqrt{n - 1}}{ \sqrt{n + 1} + \sqrt{n - 1} }
	= \lim_{n \to \infty}\frac{2 \frac{\sqrt{n - 1}}{\sqrt{n + 1}}}{ \frac{\sqrt{n + 1}}{\sqrt{n + 1}} + \frac{\sqrt{n - 1}}{\sqrt{n + 1}} } \\
	\lim_{n \to \infty} \sqrt{n^2 - 1} - n + 1
	& = \frac{2 \cdot 1}{1 + 1} = 1
\end{align*}

\subsection{סעיף ו'}
נוכיח כי $L = 1$ הוא גבול חלקי של $(a_n)$.
\begin{proof}
	נגדיר סדרת אינדקסים
	\[
		n_k = \sqrt{n^2 - 1}
	\]
	על־פי סעיף ד' מתקיים
	\[
		a_{n_k} = \left\langle \sqrt{n^2 - 1} \right\rangle = \sqrt{n^2 - 1} - n + 1
	\]
	ולכן על־פי סעיף ה'
	\[
		\lim_{k \to \infty} a_{n_k} = \lim_{n \to \infty} \sqrt{n^2 - 1} - n + 1 = 1
	\]
	לכן $L = 1$ גבול חלקי של $(a_n)$.
\end{proof}

\subsection{סעיף ז'}
נחשב את $L = \displaystyle\uplim_{n \to \infty} a_n$. \\*
נניח בשלילה כי $L > 1$, לכן קיימת תת־סדרה $a_{n_k}$ אשר מקיימת $\lim_{k \to \infty} a_{n_k} > 1$. \\*
על־פי הגדרת הגבול במצב זה לפי הגדרה 2.9 קיים $n \in \NN$ כך שמתקיים $a_n = 1$ בסתירה למסקנת סעיף א', לכן $L \le 1$.
בסעיף הקודם מצאנו כי $L = 1$ הוא גבול חלקי של הסדרה, וידוע כי $1 \ge L$ לכל גבול חלקי $L$ ולכן בהתאם הוא הגבול החלקי הגדול ביותר של $(a_n)$, דהינו
\[
	\displaystyle\uplim_{n \to \infty} a_n = 1
\]

\subsection{סעיף ח'}
נגדיר $A = \left\{ a_n \mid n \in \NN \right\}$ ונמצא את $\sup A$. \\*
אנו יודעים כי $1 \notin A$, אך על־פי הסעיף הקודם אנו יודעים כי לכל $n < 1$ בסביבה של $1$ מתקיים $n \in A$. \\*
מסיבה זו $l = 1$ הוא חסם מלעיל של $A$, ולכל $k < l$ שנבחר קיים $a \in A$ כך ש־$a > k$, לכן $l$ חסם עליון, נסמן
\[
	\sup A = 1
\]
כאמור, לקבוצה $A$ אין מקסימום, על־פי טענה $(1)$.

\end{document}
